\documentclass[12pt,letterpaper, onecolumn]{exam}
\usepackage{graphicx}
\usepackage{caption}
\usepackage{color}
\usepackage{caption}
\usepackage{listings}
\usepackage[dvipsnames]{xcolor}
\graphicspath{ {./images/} }
\lhead{Mustafa Rashid\\}
\rhead{Homework 5: Query Implementation\\}
\chead{\hline} 
\thispagestyle{empty} 
\newcommand*{\setdef}[1]{\left\{#1 \right\}} 
\renewcommand{\thepartno}{\Alph{partno}} % Force uppercase letters
\renewcommand{\partlabel}{\thepartno.}    % Format as A. B. C.

\begin{document}
\begingroup  
\noindent\LARGE CIS 5500: Database and Information Systems\\
\noindent\LARGE Homework 5: Query Implementation\\
\noindent\large \today\\
\noindent\large Mustafa Rashid\par
\noindent\large Spring 2026\par
\endgroup
\rule{\textwidth}{0.4pt}
\pointsdroppedatright
\printanswers
\renewcommand{\solutiontitle}{\noindent\textbf{Ans:}\enspace}

% SQL listing from https://gist.githubusercontent.com/vivngo/e37e1c7b79bf7e8b910c6566c59c46be/raw/db8723fc80c8b1b01665851106dba2f370b5d865/code.tex
\definecolor{dkgreen}{rgb}{0,0.6,0}
\definecolor{gray}{rgb}{0.5,0.5,0.5}
\definecolor{mauve}{rgb}{0.58,0,0.82}
\lstset{language=SQL,
  basicstyle={\small\ttfamily},
  belowskip=3mm,
  breakatwhitespace=true,
  breaklines=true,
  classoffset=0,
  columns=flexible,
  commentstyle=\color{dkgreen},
  framexleftmargin=0.25em,
  frameshape={}{}{}{}, %To remove to vertical lines on left, set `frameshape={}{}{}{}`
  keywordstyle=\color{blue},
  numbers=none, %If you want line numbers, set `numbers=left`
  numberstyle=\tiny\color{gray},
  showstringspaces=false,
  stringstyle=\color{mauve},
  tabsize=3,
  xleftmargin =1em
}

\begin{questions}
  \question \textbf{Question 1: Relational Algebra and Equivalences (15 points)}
  \begin{parts}
    \part (5 points)
    \begin{solution}
      \[\pi_{\text{SName}}(\sigma_{\text{S.SId=E.SId}\land\text{S.Age<21}\land\text{E.CId=C.CId}\land\text{C.Instructor=``Dr.Naik"}}(S\times E\times C))\]
    \end{solution}
    \part \textbf{5 points}
    \begin{solution}
      \[\pi_{\text{SName}}(\sigma_{\text{S.Age<21}\land\text{C.Instructor=``Dr.Naik"}}((S\bowtie E)\bowtie C))\]
    \end{solution}
    \part (5 points)
    \begin{solution}\\
      \begin{align*}
\pi_{\text{SName}}(&\pi_{\text{SName, CId}}(\pi_{\text{SId, SName}}(\sigma_{\text{Age<21}}S) \bowtie \pi_{\text{SId, CId}}E) \\
&\bowtie \pi_{\text{CId}}(\sigma_{\text{Instructor=``Dr.Naik"}}C))
\end{align*}
    \end{solution}
  \end{parts}
  \question \textbf{Question 2: Cost estimation (40 points)}
  \begin{parts}
    \part (5 points)
    \begin{solution}
      The Researcher relation has 8000 tuples. Out of those \(\frac{1}{50}\) are in the field of biology and \(\frac{}{}\) have been active for 12 or more years. So we 
      can estimate that \(\frac{1}{50}\dot\frac{40-12}{40}\dot8000=112\) tuples to be the result of the subquery.
    \end{solution}
    \part (10 points)
    \begin{solution}
      \begin{enumerate}
        \item The subquery \(\sigma_{\text{region=1}\lor\text{region=3}\text{Conference}}\) will have \(\frac{1}{4}\dot2000+\frac{1}{4}]dot2000=1000)\) tuples. Because each researcher is in one conference, we will have 112 tuples which we project onto two attributes. This means that we will need \(\lceil\frac{224}{400}\rceil\) or one page.
      \item The subquery \(\sigma_{\text{region=4}}\text{Conference}\) will have \(\frac{1}{4}\dot2000=500\) tuples. We then join on \texttt{cid} which will give us 500 tuples as each conference has corresponding stats. Projecting this onto 3 attributes will give us 1500 attributes meaning that we will need \(\lceil\frac{1500}{400}\rceil\) or 4 pages.
        \end{enumerate}
        \part (10 points)
        \begin{solution}
          \begin{enumerate}
            \item Because \(\lvert\text{ConfStats}\rvert=\lvert\text{Conference}\rvert\) it does not matter which relation we use as outer. The cost will be \(20+\rceil\frac{20}{18}\rceil\dot20=60\) I/Os
            \item Because \
          \end{enumerate}
        \end{solution}
    \end{solution}
  \end{parts}
  \question \textbf{Question 3:Improving queries using EXPLAIN PLAN (45 points)}
\end{questions}
\end{document}
